\chapter{The effects}\label{ch:effects}
\nwsaom{} follows RSiena's effect library. Of the \nTotal{} rows of RSiena's effects table, \nImpl{} are implemented (\pctImpl\% of the \nPlannable{} rows that are in scope), \nPartial{} is partly implemented, \nNotImpl{} are not implemented yet and \nNotPlanned{} (continuous dependent variables) are out of scope. Appendix~\ref{app:effects} lists every implemented effect with the RSiena short name, the network type and the \nwsaom{} option that selects it; Chapter~\ref{ch:missing} lists what is missing.

\section{Naming and the option families}
RSiena names an effect by a short name (\texttt{transTrip}, \texttt{egoX}, \texttt{avAlt}) and the covariate or network it interacts with. \nwsaom{} selects effects by Stata options, grouped by what the effect needs:
\begin{description}[leftmargin=0pt,labelindent=0pt]
\item[Structural effects] switches such as \opt{outdegree}, \opt{reciprocity}, \opt{transtrip}, \opt{cycle3}, \opt{indegpopularity}, \opt{outactivity}, \opt{gwesp}; effects with a parameter use \opt{outtrunc(\#)}, \opt{gwespfb(\#)} etc., and the additional own-network effects of the later versions are collected in \opt{ownx(effect[:\#] \ldots)}.
\item[Covariate effects] \opt{nodeocov()}/\opt{egox()}, \opt{nodeicov()}/\opt{altx()}, \opt{nodematch()}/\opt{samex()}, \opt{simcov()}/\opt{simx()}, \opt{egosqx()}, \opt{altsqx()}, \opt{absdiffx()}, \ldots, and \opt{netx(effect:var)} for the further covariate effects. Each takes a variable list. A changing covariate (RSiena \texttt{varCovar}) is declared with \opt{tvcov()} and then used by its stem name.
\item[Dyadic covariates] \opt{dyadcov()}, \opt{dyadcovrecip()}, \opt{dyadx(effect:matrix)}, and with changing matrices \opt{tvdyadcov()}, \opt{tvdyadcovrecip()}, \opt{tvdyadx()}.
\item[Behavior effects] \opt{behavior(vars : shape effects)} with \opt{beheff()}, \opt{behx()}, \opt{behw()}, \opt{tvbehw()}, \opt{behshape()}, \opt{behnet2()}; and the selection effects of the dependent behavior on the network: \opt{behsim}, \opt{behegox}, \opt{behaltx}, \opt{netbeh()}.
\item[Several networks] \texttt{nwsaom multinet} with \opt{cross1()}/\opt{cross2()} (Chapter~\ref{ch:networktypes}).
\item[Rate effects] \opt{ratecov()}, \opt{outrate}, \opt{inrate}, \ldots (Chapter~\ref{ch:rates}).
\end{description}
Where RSiena offers a parameter (the RSiena ``internal effect parameter'') the option takes it after a colon, for example \opt{ownx(outinass:2)} or \opt{behnet2(totgwdspffalt:69)}; the defaults are RSiena's.

\section{Statistics, contributions and what ``validated'' means}
For every effect RSiena defines a \emph{contribution} (used when the simulation evaluates a ministep) and a \emph{tie or ego statistic} (used for the target of the method of moments). \nwsaom{} implements both from the RSiena source, including the cases in which the two differ, and checks them in two ways: the observed statistic is compared with the target that \texttt{siena07} reports for the same data, and the estimates are compared over at least two seeds (the estimates of weakly identified effects agree only within their standard errors). The column ``Validated'' in Appendix~\ref{app:effects} says \texttt{siena07} when this comparison was made. Every implemented effect has been compared this way; for a few effects the comparison rests on the statistic alone, because the estimates are degenerate or RSiena returns noise for them (the notes in Appendix~\ref{app:effects} say which). Where RSiena's own behaviour is odd (it counts a tie twice, reads a covariate at another node's index, \ldots) the quirk is copied on purpose (Appendix~\ref{app:quirks}).

\section{Effect families added in the later versions}
Besides the effects of the single directed network, the library now covers the following families; each is described where its network type is treated.
\begin{itemize}[nosep]
\item \textbf{Rate effects} (Chapter~\ref{ch:rates}): the structural rate effects \opt{outrate}, \opt{outrateinv} and \opt{outratelog} for non-directed and two-mode networks; \opt{ratecov()} for two-mode networks, with changing covariates and with the behavior as the covariate; and the rate effects of the behavior in a co-evolution model (\opt{behratecov()}, \opt{behrate()}).
\item \textbf{Two-mode networks} (Section~\ref{sec:twomode}): \opt{dyadcov()} with an actors-by-events matrix, \opt{interact()}, and \opt{netbeh(crossxoutact avgroupegox totgroupegox)} for the behavior of the actors.
\item \textbf{Several networks} (Section~\ref{sec:multinet}): a two-mode network as one of the two, the 21 two-mode cross effects of RSiena and their counterparts for a one-mode network that depends on a two-mode one, \texttt{toany}, \texttt{gwespmix}, \texttt{betweenpop}, \texttt{from.w.ind}, the dyadic-covariate cross effects \texttt{toU}, \texttt{sharedToU} and \texttt{avAltU}, changing covariates in the cross covariate effects, and cross effects for non-directed pairs; with a behavior the behavior effects of both networks, covariate, shape and network effects of the behavior, the effects \texttt{netbeh:}, endowment and creation of behavior effects, missing behavior values and interactions of behavior effects.
\item \textbf{Effects of RSiena 1.6.14}: \texttt{toU} (whose two networks were exchanged in 1.6.12), \texttt{toU} and \texttt{WXX} for two-mode networks, and \texttt{WXZX}.
\end{itemize}
Effects that RSiena cannot evaluate are not offered (Section~\ref{sec:cannoteval}).

\section{Effects that are native only}
Not every effect can be simulated without the compiled plugin; Appendix~\ref{app:mata} lists the effects that also run in Mata. All new effects, in particular those of the families above, exist only in the plugin, which is part of the package. For all of these \nwsaom{} stops with an explicit message if the plugin is missing or too old.
